\documentclass[10pt, aspectratio=169]{beamer}
\input{../../_en_preambulo_beamer}
% Comandos y macros estadísticas compartidas para las presentaciones Beamer en inglés (Metropolis)

% Conjuntos numéricos básicos
\newcommand{\R}{\mathbb{R}}
\newcommand{\N}{\mathbb{N}}
\newcommand{\Q}{\mathbb{Q}}
\newcommand{\Z}{\mathbb{Z}}
\newcommand{\set}[1]{\left\{ #1 \right\}}
\newcommand{\abs}[1]{\left|#1\right|}
\newcommand{\norm}[1]{\left\| #1 \right\|}

% Comandos estadísticos y probabilísticos del curso
\newcommand{\Var}{\operatorname{Var}}
\newcommand{\cov}{\operatorname{Cov}}
\newcommand{\corr}{\rho}
\newcommand{\comb}[2]{\begin{pmatrix} #1 \\ #2 \end{pmatrix}}
\newcommand{\s}{\sigma}
\newcommand{\particion}{\mathfrak{P}}
\newcommand{\card}[1]{\#\left( #1 \right)}
\newcommand{\seq}[3]{\set{#1}_{#2}^{#3}}
\newcommand{\E}{\mathbb{E}}
\newcommand{\Prob}{\mathbb{P}}

% Entornos pedagógicos y cajas en inglés académico natural (soportando tanto nombres en inglés como en español)
\newenvironment{discussion}{%
  \begin{alertblock}{Discussion Question}%
}{%
  \end{alertblock}%
}
\newenvironment{discusion}{%
  \begin{alertblock}{Discussion Question}%
}{%
  \end{alertblock}%
}

\newenvironment{reflection}{%
  \begin{alertblock}{Classroom Reflection}%
}{%
  \end{alertblock}%
}
\newenvironment{reflexion}{%
  \begin{alertblock}{Classroom Reflection}%
}{%
  \end{alertblock}%
}

\newenvironment{paradox}{%
  \begin{alertblock}{Exploratory Paradox}%
}{%
  \end{alertblock}%
}
\newenvironment{paradoja}{%
  \begin{alertblock}{Exploratory Paradox}%
}{%
  \end{alertblock}%
}

\newenvironment{motivation}{%
  \begin{exampleblock}{Motivation: Data Science in Practice}%
}{%
  \end{exampleblock}%
}
\newenvironment{motivacion}{%
  \begin{exampleblock}{Motivation: Data Science in Practice}%
}{%
  \end{exampleblock}%
}

\newenvironment{quickchallenge}{%
  \begin{block}{Quick Team Challenge}%
}{%
  \end{block}%
}
\newenvironment{retoclase}{%
  \begin{block}{Quick Team Challenge}%
}{%
  \end{block}%
}


\title[Mathematics of Regression]{Section 09.03: The Mathematics of Linear Regression}
\subtitle{OLS Derivation, the SST=SSR+SSE Decomposition, and $R^2$ Properties}
\author[J. Castillo Colmenares]{Juliho Castillo Colmenares}
\institute{Tecnológico de Monterrey}
\date{\vspace{-1.5cm}}

\begin{document}

% Slide 1: Title
\begin{frame}[plain]
  \titlepage
\end{frame}

% Slide 2: Roadmap
\begin{frame}{Roadmap --- Chapter 09: Linear and Multiple Regressions}
  \scriptsize
  Where are we in the modular development of this chapter?
  \begin{itemize}\itemsep=0.02em
    \item \textbf{09.01} Correlation as the Premise of Regression
    \item \textbf{09.02} Introduction to Linear Regression
    \item \textbf{09.03} The Mathematics of Regression: OLS Derivation and $R^2$ \emph{(Today)}
    \item \textbf{09.04} Linear Regression on Simulated Data
    \item \textbf{09.05} Optimal Coefficients, $t$/$F$ Tests, and RSE
    \item \textbf{09.06} Implementation in Python with \texttt{statsmodels}
    \item \textbf{09.07} Multiple Regression, the Normal Equation, and Ridge/Lasso
    \item \textbf{09.08} Model Validation and $k$-fold Cross-Validation
    \item \textbf{09.09} Residual Diagnostics and Multicollinearity (VIF)
    \item \textbf{09.10} Regression with \texttt{scikit-learn}
    \item \textbf{09.11} Categorical Variables and Dummy Variables
    \item \textbf{09.12} Nonlinear Transformations and Polynomial Regression
  \end{itemize}
  \pause
  \vspace{-0.05cm}
  \begin{block}{Session Objective}
    Rigorously derive the Ordinary Least Squares (OLS) formulas via differential calculus, prove the fundamental decomposition $\text{SST}=\text{SSR}+\text{SSE}$, and establish the properties of the coefficient of determination $R^2$.
  \end{block}
\end{frame}

% Slide 3: Motivation
\begin{frame}{The Goal: Minimizing the Fitting Error}
  \footnotesize
  \begin{motivacion}
    We propose the model $Y_e=\hat\alpha+\hat\beta X$ to estimate $Y$ from $X$. The purpose of linear regression is to find the values $\hat\alpha,\hat\beta$ that \emph{minimize} the difference between $Y$ and $Y_e$, represented by the error $\epsilon=Y-Y_e$.
  \end{motivacion}

  \vspace{0.15cm}
  \pause
  The \textbf{least squares} method minimizes the sum of squared errors, not absolute errors: this makes the problem differentiable and penalizes large errors more heavily.
\end{frame}

% Slide 4: Objective function
\begin{frame}{The Least Squares Objective Function}
  \footnotesize
  \begin{block}{Objective Function}
    $$S(\alpha,\beta) = \sum_{i=1}^n\left(Y_i-(\alpha+\beta X_i)\right)^2 \quad \longrightarrow \quad \min_{\alpha,\beta} S$$
  \end{block}
  \pause

  \vspace{0.1cm}
  \begin{alertblock}{Minimization Strategy}
    We set the partial derivatives $\partial S/\partial\alpha=0$ and $\partial S/\partial\beta=0$ to zero: the resulting system is known as the \textbf{normal equations}.
  \end{alertblock}
\end{frame}

% Slide 5: Derivation
\begin{frame}{Deriving the Least Squares Formulas}
  \footnotesize
  \begin{block}{Derivative with respect to $\alpha$}
    $$\frac{\partial S}{\partial\alpha}=-2\sum_i(Y_i-\alpha-\beta X_i)=0 \implies \hat\alpha=\bar Y-\hat\beta\bar X$$
  \end{block}
  \pause

  \vspace{0.1cm}
  \begin{alertblock}{Derivative with respect to $\beta$ (substituting $\hat\alpha$)}
    $$\hat\beta=\frac{\sum_i(X_i-\bar X)(Y_i-\bar Y)}{\sum_i(X_i-\bar X)^2}=\frac{S_{XY}}{S_{XX}}$$
  \end{alertblock}
\end{frame}

% Slide 6: Partition theorem
\begin{frame}{Partition Theorem: SST = SSR + SSE}
  \footnotesize
  \begin{block}{Partition Identity}
    $$\underbrace{\sum(Y_i-\bar Y)^2}_{\text{SST}} = \underbrace{\sum(\hat Y_i-\bar Y)^2}_{\text{SSR}} + \underbrace{\sum(Y_i-\hat Y_i)^2}_{\text{SSE}}$$
  \end{block}
  \pause

  \vspace{0.1cm}
  \begin{alertblock}{Key to the Proof}
    The cross term $\sum(\hat Y_i-\bar Y)(Y_i-\hat Y_i)$ vanishes exactly because the normal equations guarantee $\sum e_i=0$ and $\sum X_ie_i=0$ (and hence $\sum\hat Y_ie_i=0$).
  \end{alertblock}
\end{frame}

% Slide 7: R²
\begin{frame}{The Coefficient of Determination $R^2$}
  \footnotesize
  \begin{block}{Definition and Properties}
    $$R^2=\frac{\text{SSR}}{\text{SST}}=1-\frac{\text{SSE}}{\text{SST}} \in[0,1], \qquad R^2=1 \iff \text{SSE}=0$$
  \end{block}
  \pause

  \vspace{0.1cm}
  \begin{alertblock}{Adjusted $R^2$}
    $$R^2_{\text{adj}}=1-\frac{n-1}{n-p-1}(1-R^2)$$
    $R^2$ never decreases when adding predictors; $R^2_{\text{adj}}$ can decrease if the new predictor doesn't contribute enough explanatory power.
  \end{alertblock}
\end{frame}

% Slide 8: Python Lab (Block 1)
\begin{frame}[fragile]{Python Lab: Verifying the Normal Equations}
  \scriptsize
  Numeric confirmation that OLS satisfies $\sum e_i=0$ and $\sum X_ie_i=0$ (\texttt{09.03\_mathematics\_of\_regression.py}):
  {\lstset{basicstyle=\fontsize{3.6pt}{4.3pt}\selectfont\ttfamily,aboveskip=0.05em,belowskip=0.05em}
  \lstinputlisting[firstline=17, lastline=27]{../../code/09_regresiones/09.03_mathematics_of_regression.py}}
\end{frame}

% Slide 9: Python Lab (Block 2)
\begin{frame}[fragile]{Python Lab: SST=SSR+SSE Decomposition}
  \scriptsize
  Verifying the identity and the vanishing cross term that proves it:
  {\lstset{basicstyle=\fontsize{3.4pt}{4.1pt}\selectfont\ttfamily,aboveskip=0.05em,belowskip=0.05em}
  \lstinputlisting[firstline=32, lastline=44]{../../code/09_regresiones/09.03_mathematics_of_regression.py}}
\end{frame}

% Slide 10: Python Lab (Block 3)
\begin{frame}[fragile]{Python Lab: $R^2$ Properties}
  \scriptsize
  The $R^2=r^2$ identity and adjusted $R^2$ behavior facing a useless predictor:
  {\lstset{basicstyle=\fontsize{3.2pt}{3.9pt}\selectfont\ttfamily,aboveskip=0.05em,belowskip=0.05em}
  \lstinputlisting[firstline=48, lastline=69]{../../code/09_regresiones/09.03_mathematics_of_regression.py}}
\end{frame}

% Slide 11: Terminal Output
\begin{frame}[fragile]{Python Lab: Terminal Output}
  \scriptsize
  Standard output produced when running the lab script:
  \vspace{0.02cm}
  \begin{block}{Standard Console Output}
    \fontsize{3.6pt}{4.3pt}\selectfont
    \begin{verbatim}
=== Block 1: Normal Equations ===
sum(residuals)=0.0000000000, sum(x*residuals)=0.0000000000

=== Block 2: SST=SSR+SSE ===
SST=14309.41, SSR=12304.00, SSE=2005.41; cross term=0.0000000000

=== Block 3: R^2 Properties ===
R^2=0.8599 = r^2=0.8599
Adding useless predictor: R^2 0.8599->0.8602 (up); adj-R^2 0.8591->0.8587 (down)
    \end{verbatim}
  \end{block}
\end{frame}


% Slide 15: Closing
\begin{frame}{Section Closing: Mathematics of Regression}
  \begin{reflexion}
    We have rigorously derived the OLS formulas from first principles (differential calculus), proven that the SST=SSR+SSE decomposition is an exact consequence of the normal equations, and established the fundamental properties of the $R^2$ coefficient.
  \end{reflexion}

  \vspace{0.2cm}
  \pause
  \begin{block}{Key Takeaways}
    \begin{itemize}\itemsep=0.08cm
      \item \textbf{OLS:} $\hat\beta=S_{XY}/S_{XX}$, $\hat\alpha=\bar Y-\hat\beta\bar X$, unbiased by construction.
      \item \textbf{Orthogonality:} $\sum e_i=0$ and $\sum X_ie_i=0$ are the algebraic key to all subsequent theory.
      \item \textbf{$R^2$ vs. $R^2_{\text{adj}}$:} the former never decreases; the latter penalizes unhelpful predictors.
    \end{itemize}
  \end{block}
\end{frame}

% Slide 16: Modular Perspective
\begin{frame}{Modular Perspective --- The Next Challenge}
  \scriptsize
  \begin{retoclase}
    With the formulas rigorously derived, Section 09.04 applies them to a simulated dataset generated explicitly with known parameters (true $\alpha,\beta,\sigma$), allowing a direct comparison between the OLS-estimated values and the actual population values that generated the data.
  \end{retoclase}

  \vspace{0.08cm}
  \pause
  \begin{alertblock}{Recommended Preparation}
    \begin{itemize}\itemsep=0.06cm
      \item Review the difference between population parameters ($\alpha,\beta$) and their sample estimators ($\hat\alpha,\hat\beta$).
      \item Reflect on how close you'd expect $\hat\beta$ to be to $\beta$ as the sample size $n$ grows.
    \end{itemize}
  \end{alertblock}
\end{frame}

\end{document}
